\section{Summing the trace}\label{sec:trace-count}

We now prove the high-trace estimate.  The input from
\Cref{prop:retained-lines} is a description of the activities of the fixed
primes on the first-visit tree.  The input from
\Cref{prop:cylinder-reduction} replaces vertex deletion by short lists of
qualifying primitive specifications.  The remaining task is to sum the
resulting conditional averages without losing a constant for every prime
occurrence.  An exact identity for weighted subtrees makes this possible.

We retain the notation $\mathscr T$, $b(v)$ and $r_*$ for the first-visit
tree, its integer offsets and the number of returns.  In this section a
subtree is a connected subgraph of $\mathscr T$; its \emph{top} is its
vertex nearest the root.  Put
\[
 s(v)=\#\{\text{children of }v\},\qquad
 M_0=t_0(LJ+1),\qquad N_{\mathrm{exc}}=B^{1-\rho/2}.
\]
We call a tree edge $u\longrightarrow v$ \emph{good} if
$s(u)=s(v)=1$, $u$ is not the root, and neither endpoint belongs to any
return step.  Write $\mathcal G$ for the set of good edges.  These
conditions ensure that an uncorrupted prime occurring just on such an
edge has exactly two normalization factors on its lit residue.

For all but a small exceptional set of good edges, we will obtain the
factor $B^{-1+\eta}$ from a core-prime restriction to the divisor interval,
or a stronger penalty from a core cutoff.  Returns receive a cutoff
penalty as well.  Signed center-prime averages provide the further saving
needed in the high-trace estimate.  The weighted-subtree identity below
will let us sum the prime labels without losing these gains.

\subsection{Tree geometry and weighted subtrees}

\begin{lemma}\label{lem:trace-topology}
For a closed line of length $\ell$ with $r_*$ returns, the topology of first
visits and returns, together with all step signs, has at most
$4^\ell(\ell+1)^{r_*}$ possibilities.  Its tree satisfies
\begin{equation}\label{eq:trace-bad-edges}
 |E(\mathscr T)\setminus\mathcal G|\le 10(r_*+1),\qquad
 \ell=|E(\mathscr T)|+r_*.
\end{equation}
Once the labels of the tree edges and this topology are specified, every
return label is determined.
\end{lemma}

\begin{proof}
At each step record whether its destination is new and record its sign.
At each return record one of at most $\ell+1$ previously visited vertices.
These data give the stated bound.  A tree-edge label and its sign specify
the difference of its endpoint offsets; all offsets are consequently
determined, and a return label must be the absolute endpoint difference
divided by $h$.  We retain only assignments for which the abstract vertices
have distinct offsets and all these return labels belong to $\mathcal D$.
There is no independent choice of a return label.

Let $q$ be the number of leaves and $b$ the number of vertices with at
least two children.  The adding steps occur in at most $r_*+1$ runs,
each a downward path, so $q\le r_*+1$.  Also $b\le q-1$ and
\[
 \sum_{s(v)\ge2}s(v)=q+b-1.
\]
The number of edges incident with a leaf or such a branch vertex is
therefore at most $4q$.  There are at most $2r_*$ endpoints of return
steps.  An endpoint that has not already been counted is unary and is
incident with at most two tree edges, giving at most $4r_*$ further
edges.  The root adds at most one edge unless it was already counted as a
branch vertex.  These bounds imply \eqref{eq:trace-bad-edges}.
\end{proof}

For $a\in\{C,Z\}$ and a subtree $Q$ with at least one edge define
\begin{equation}\label{eq:subtree-weights}
 u_a(Q)=A_a^{|E(Q)|}\prod_{v\in V(Q)}\beta_a^{s(v)},\qquad
 S_a(\mathscr T)=\sum_{v\in V(\mathscr T)}(1-\beta_a^{s(v)}).
\end{equation}
The degrees in this definition are those of the full tree.

\begin{lemma}\label{lem:subtree-weight}
For every finite rooted tree and every $A_a>0$ with
$\beta_a=(1+A_a)^{-1}$,
\begin{equation}\label{eq:subtree-identity}
 \sum_{Q:\,|E(Q)|\ge1}u_a(Q)=S_a(\mathscr T)\le |V(\mathscr T)|.
\end{equation}
\end{lemma}

\begin{proof}
Start at a prescribed vertex $v$.  At every reached vertex include each
child edge independently with probability $A_a\beta_a$ and exclude it
with probability $\beta_a$.  The resulting connected subtree has top
$v$.  A particular outcome $Q$ has probability
\[
 (A_a\beta_a)^{|E(Q)|}
 \beta_a^{\sum_{w\in V(Q)}s(w)-|E(Q)|}=u_a(Q).
\]
The outcome with no edges has probability $\beta_a^{s(v)}$.
Sum the other probabilities, then sum over the possible tops $v$.
\end{proof}

The exact sum in \eqref{eq:subtree-identity} will cancel the background
exponential in the primewise estimate below.  To retain the additional
center-prime saving, we also use modified weights.  These anticipate the
primewise bounds proved in \Cref{lem:trace-integration}: the reduction on
a single good edge comes from a signed average, while the increase on
other single edges permits an absolute bound.
Set $\widetilde u_C=u_C$.  In the center
band keep $\widetilde u_Z(Q)=u_Z(Q)$ unless $Q$ consists of one edge;
for a single good edge set $\widetilde u_Z(Q)=\theta/2$, and for a single
nongood edge set $\widetilde u_Z(Q)=u_Z(Q)+\theta$.  Since both endpoint
degrees of a good edge are one, its original weight is
$\beta_Z^2=\theta$.  Thus
\begin{equation}\label{eq:modified-subtree-sum}
 \sum_Q\widetilde u_Z(Q)-S_Z(\mathscr T)
 =-\frac\theta2|\mathcal G|
       +\theta|E(\mathscr T)\setminus\mathcal G|.
\end{equation}

\subsection{Recording prime labels}

Fix a retained line and a list $\mathfrak L$ of fewer than $t_0$
primitive specifications attached at its vertices, as supplied by
\Cref{prop:cylinder-reduction}.  Write $I_{\mathfrak L}$ for their joint
qualification indicator.  Recall that a \emph{tree prime} occurs in
a tree-edge label.  As before, the fixed label set $\mathcal F$ consists
of the core primes occurring on the line and all primes occurring at
least twice on the line.  A fixed prime is active at $v$ when
$p\mid n+b(v)$.  Active components include isolated vertices.

Tag a tree prime if it satisfies at least one of the following conditions:
\begin{enumerate}[label=(\roman*),nosep]
 \item it occurs among the primes of $\mathfrak L$;
 \item it is fixed and has several active components, or has an unlit
 occurrence anywhere on the line;
 \item it is nonfixed and corrupted in the sense of
 \Cref{prop:retained-lines}.
\end{enumerate}
These decisions depend only on the line, the list and the fixed
coordinates $n_{\mathcal F}$.

We record tree labels by the following finite data, called \emph{tokens}.
For a tagged fixed prime, record one subtree token for each active
component having an edge, and one single-edge token for each unlit tree
occurrence.  The latter is called a ghost token and is assigned weight
$2$.  A subtree token of band $a$ and shape $Q$ has weight $u_a(Q)$.
For a tagged nonfixed prime, record its sole tree occurrence by a ghost
token.  Tokens belonging to the same tagged prime are grouped together,
and the group is assigned that prime value.

Every untagged tree prime receives a single token.  If it is fixed, its
active set is connected and all its occurrences are lit; the token is
the subtree formed by its tree occurrences.  If it is nonfixed, its token
is its single tree edge.  An untagged token of band $a$ and shape $Q$
has weight $\widetilde u_a(Q)$.  List prime slots, including the extra
modulus prime of each specification, are recorded along with their
equalities to each other and to tagged groups.  An untagged prime occurs
in no such group or slot.

\begin{lemma}\label{lem:trace-token-record}
For retained data, these records have the following properties.
\begin{enumerate}[label=(\roman*)]
 \item They recover every tree-edge label.  Together with the topology
 and the list data they consequently recover the full line and list.
 \item The total number of tagged tokens and list prime slots is at
 most $N_{\mathrm{exc}}$, for sufficiently large $B$.
 \item An untagged fixed center prime cannot have a single good edge
 as its token.
\end{enumerate}
\end{lemma}

\begin{proof}
If adjacent vertices are active for $p$, their difference $\pm hd$ is
divisible by $p$.  Since $p\nmid h$, the intervening label contains $p$.
Conversely, a lit occurrence has both endpoints active.  Thus the
nontrivial active components record exactly the lit tree occurrences.
Ghosts record the other occurrences; isolated active vertices create no
label.  Multiplying the distinct prime values of the tokens covering an
edge recovers its label, proving (i).

By \Cref{prop:retained-lines}, the fixed primes with several active
components contribute $O(B^{1-2\rho}B^\rho)=O(B^{1-\rho})$ subtree
tokens.  Unlit fixed occurrences contribute $O(t_0)$ ghosts and at most
that many newly tagged connected labels.  Corrupted nonfixed primes
contribute $O(B^{1-2\rho})$ ghosts.  Lists contribute $O(M_0)$ slots
and at most one new token per newly tagged connected tree prime:
disconnected primes and primes with unlit occurrences have already been
counted.  Since
\[
 M_0=O(B^{1-\rho+4\eps}\log B),\qquad
 1-\rho+4\eps<1-\rho/2,
\]
the sum of these bounds is at most $N_{\mathrm{exc}}$ eventually.

For (iii), such a prime is repeated on the line and all its occurrences
are lit.  If its tree token consists of a single edge, its connected
active set contains only the two endpoints of that edge.  A further
occurrence cannot be another tree edge, and hence must be a return
touching these endpoints.  This contradicts goodness.  An active return
endpoint elsewhere would give another active component and would have
tagged the prime.
\end{proof}

We will use the divisor interval at a good edge by restricting the prime
value of an untagged core token whose top is its origin.  First set aside
the tagged core tops: skip every good edge whose origin is the top of a
tagged core subtree token.  At most $N_{\mathrm{exc}}$ edges are skipped.
Among the remaining good edges, let $\mathcal U$ be the set of origins
at which no untagged core token has its top.  Thus no core subtree token
has top $u\in\mathcal U$, and $s(u)=1$.  At these origins we will use
a cutoff penalty instead.  The next estimate retains this penalty;
the subsequent sum over prime assignments will use the divisor interval
at the other non-skipped good origins.

For a fixed numerical line and list, group the retained fixed-coordinate
configurations according to their complete assigned token records.
Within a group the token shapes, prime values and groupings are fixed;
unrecorded isolated activities may still vary.  Order components and
tokens by any deterministic rule, so that these groups form a
partition, rather than a choice of representations for each residue
configuration.  We next bound the integral of each such group.  This is
the point at which the normalization of the graph produces cancellation.

\subsection{Primewise integration}

For an assigned record $R$, let $\1_R(n_{\mathcal F})$ denote its group
indicator.  Its contribution to the conditional expression is
\[
 \mathcal I(R)=
 \E_{\mathcal F}\left[
  \1_R\left|
   \E_{\mathcal S\setminus\mathcal F}
       K_{\boldsymbol d}(n)I_{\mathfrak L}(n)
       \right|\right].
\]
Empty or incompatible groups have contribution zero.  Products over
tagged tokens below use the weights just defined, including weight $2$
for ghosts.

\begin{lemma}\label{lem:trace-integration}
Uniformly over retained lines, lists and assigned records,
\begin{align}
 \mathcal I(R)
 &\le
 \exp\left\{
 o(1)-\kappa(v_C-T\sqrt{v_C})(r_*+|\mathcal U|)\right\}
 \notag\\
 &\quad\times\exp\left\{
 -v_CS_C(\mathscr T)-v_ZS_Z(\mathscr T)
 +v_C\beta_C(e^\kappa-1)|\mathcal U|\right\}
 \notag\\
 &\quad\times\prod_{\substack{p\text{ distinct}\\
                  p\text{ in tree or list}}}\frac1p
 \prod_{\text{tagged tokens}}\mathrm{weight}
 \prod_{\substack{\text{untagged tokens}\\(a,Q)}}\widetilde u_a(Q).
 \label{eq:trace-integration}
\end{align}
Primes occurring only on returns need not appear in the product over
tree or list primes.  The $o(1)$ tends to zero as $B$ tends to infinity,
uniformly in the line, list and record, with $h,\tau,C_0,T$ and the
constants in \eqref{eq:graph-constants} fixed.
\end{lemma}

\begin{proof}
\emph{Order of integration.}
The set $\mathcal F$ does not contain every core prime.  Put
$\mathcal C'=\mathcal C\setminus\mathcal F$ and
$\mathcal Z'=\mathcal Z\setminus\mathcal F$.  At fixed
$n_{\mathcal F}$ use exactly the inequality
\begin{equation}\label{eq:trace-fubini}
 \left|\E_{\mathcal C'}\E_{\mathcal Z'}
             K_{\boldsymbol d}I_{\mathfrak L}\right|
 \le \E_{\mathcal C'}
       \left|\E_{\mathcal Z'}
             K_{\boldsymbol d}I_{\mathfrak L}\right|.
\end{equation}
In the inner average hold \emph{all} core coordinates as parameters.
The cutoffs can depend on the background core coordinates in
$\mathcal C'$, but never on $\mathcal Z'$.  Apart from these cutoffs,
the kernel and each compatible list cylinder separate by prime.
The inner center-prime averages therefore factor, with their signs
intact.  We will subsequently use positive majorants in the other
coordinates.  No enlargement of the fixed set in the cylinder
construction is involved.

\emph{Charging the cutoffs.}
Use the origin cutoff at each return step.  For each $u\in\mathcal U$,
use the endpoint cutoff at $u$ of its incoming tree edge, which exists
because a good origin is not the root.  These are distinct cutoff
occurrences: vertices of $\mathcal U$ touch no return, and an incoming
tree edge has a unique child.  A charged cutoff with label $d$ is at
most
\[
 \exp\{-\kappa(v_C-T\sqrt{v_C})\}
 \exp\left\{\kappa
    \sum_{\substack{p\in\mathcal C\\p\nmid d}}
           \1_{p\mid n+b(u)}\right\}.
\]
This inequality holds also off its support.  The other cutoff factors
are at most one.  The constant parts give the first negative term in
the exponent of \eqref{eq:trace-integration}.

For a return, allocate the normalization at its origin to that return.
For a core prime its multiplier is $A\beta_C$ when the prime is in the
label and lit, zero when it is in the label and unlit, and at most
$\max(1,\beta_Ce^\kappa)=1$ otherwise.  Both
$A\beta_C<1$ and $\beta_Ce^\kappa<1$ follow from $A=e^{2\kappa}$.
Center-prime return multipliers also have absolute value at most one.
These bounds will be used except when integrating a good singleton,
where we retain the return factors until its signed average has been
taken.

\emph{Fixed tree primes.}
After removing the bounded return factors, the core tree factors for
an active component $Q$ containing an edge are
\[
 A^{|E(Q)|}\prod_{v\in V(Q)}\beta_C^{s(v)}=u_C(Q).
\]
At a charged vertex $u\in\mathcal U$, activity in such a component
either continues through the incoming label, in which case that prime
is excluded from this charge, or begins a subtree token at $u$, which
is excluded by the definition of $\mathcal U$.  Thus an additional
$e^\kappa$ can occur only at an isolated active vertex there.  Its
normalization is $\beta_C^{s(u)}e^\kappa=\beta_Ce^\kappa\le1$.
Other isolated activities also contribute at most one.  The remaining
bound is the product of the core subtree weights.

A core tree prime has a lit occurrence.  Choose one deterministically
from its tokens and retain its activity indicator.  This is one
specified residue modulo $p$ and supplies the factor $1/p$ upon
integration.  For a fixed center tree prime, the same argument gives
the product of its active-subtree weights; lit factors $1-\theta/p$
are at most one, and ghost occurrences are bounded absolutely.  If
there is a lit tree occurrence, retain one such residue indicator.
If there is none, at least one unlit tree occurrence supplies
$\theta/p\le1/p$ directly, without a residue restriction.  Ghost
weight $2$ bounds all the remaining ghost factors.  These constructions
are positive primewise majorants on the assigned group.

\emph{Nonfixed tree primes.}
A tagged nonfixed tree prime has one occurrence.  Its absolute
integral is at most $2/p$: the lit residue contributes at most $1/p$
and the other residues at most $\theta/p$.  If the list fixes its
residue, the bound only improves.  An untagged nonfixed prime on a
nongood edge $Q$ contributes at most
\[
 \frac{u_Z(Q)}p+\frac\theta p
       =\frac{\widetilde u_Z(Q)}p.
\]
Indeed its lit residue includes the normalizations of both endpoints,
and any additional activities can only reduce their product.

Now let $p$ be an untagged nonfixed prime on a good edge $Q=(u,v)$.
It occurs nowhere else on the line and in no list.  On its lit residue,
noncorruption says that no other tree vertex is active.  Each endpoint
has exactly one tree departure and no return incidence.  Its entire
prime multiplier, including return factors, is therefore exactly
\[
 (1-\theta/p)\beta_Z^2=\theta(1-\theta/p).
\]
Let $a_0$ be its lit residue, and for $a\in\Z/p\Z$ define
\[
 \nu_p(a)=\#\{0\le i<\ell:a+b_i\equiv0\pmod p\}.
\]
This counts departures with multiplicity, including return departures.
We have $\nu_p(a_0)=2$.  On every other residue the entire prime
multiplier is exactly $-\theta\beta_Z^{\nu_p(a)}/p$.  Consequently its
signed average is exactly
\begin{align*}
 &\frac\theta p(1-\theta/p)
  -\frac\theta{p^2}\sum_{a\ne a_0}\beta_Z^{\nu_p(a)}\\
 &\qquad=\frac\theta{p^2}\left(
  1-\theta+
  \sum_{\substack{a\ne a_0\\\nu_p(a)>0}}
       (1-\beta_Z^{\nu_p(a)})\right).
\end{align*}
There are at most $r_p-1$ terms in the last sum, where
$r_p\le\ell+1$ is the number of residue classes represented by tree
vertices.  The displayed quantity is nonnegative and at most
\[
 \frac{\theta r_p}{p^2}
 \le\frac{\theta(\ell+1)}{p^2}
 \le\frac{\theta}{2p}
 =\frac{\widetilde u_Z(Q)}p
\]
for sufficiently large $B$.  Thus $\theta=\beta_Z^2$ cancels the
terms of order $1/p$ before any absolute value is taken.
Removing return factors by an absolute estimate before taking this
average would not justify the calculation on the no-activity residues;
here they have been retained throughout.

\emph{List primes and background primes.}
For a list prime absent from the tree, qualification supplies one
specified residue.  All other factors, including charged factors,
are bounded by one: any activity in the tree is isolated, and at a
vertex of $\mathcal U$ its factor is again $\beta_Ce^\kappa\le1$.
This yields $1/p$ for that prime.

For any remaining prime, discard return factors by the absolute bounds
already proved.  The positive tree multiplier is
\[
 \prod_{v\text{ active}}\beta_a^{s(v)}
 \exp\{\kappa\1_{a=C}\#(\mathcal U\cap\{v\text{ active}\})\}
 \le1.
\]
If the tree offsets are pairwise distinct modulo $p$, its integral is
exactly
\begin{equation}\label{eq:trace-background}
 1-\frac{S_a(\mathscr T)
       -\1_{a=C}\beta_C(e^\kappa-1)|\mathcal U|}{p}.
\end{equation}
For offset collisions we use the bound one.  A prime occurring only on
a return divides the nonzero difference between that return's two
distinct vertices, so is automatically one of these collision primes.
It entails no independently chosen prime slot.

The omitted background primes are tree or list primes, or divide one
of the $O(\ell^2)$ nonzero differences of size $O(\ell H)$.
There are $O(B\log B)$ tree and list slots.  Since every prime is
at least $P_0$ and $\log(\ell H+2)=O(B+\log B)$, their reciprocal
mass, multiplied by $\ell+1$, is $o(1)$, uniformly in the valid
assignment.  For example each nonzero difference has at most
$\log(O(\ell H)+2)/\log P_0$ prime divisors from $\mathcal S$.
Using $1-x\le e^{-x}$ in \eqref{eq:trace-background} and restoring
the omitted harmonic masses gives the remaining exponent in
\eqref{eq:trace-integration}, with an $o(1)$ error.  The numerator in
\eqref{eq:trace-background} is nonnegative, since at every charged
vertex $\beta_C e^\kappa\le1$.

Finally, the grouping does not assert independence of conditioned
coordinates.  After \eqref{eq:trace-fubini} and the signed free-center
integrations, use the positive majorants above on the group.  Whenever
a used-prime probability $1/p$ is claimed, its chosen activity or list
residue indicator is retained.  Drop all other group restrictions and
integrate the resulting product over the remaining coordinates.  This
is an enlargement of a nonnegative integral and hence is valid despite
any dependencies within the original group.  It yields one bound per
assigned record, not one bound for each individual fixed-coordinate
configuration.  Combining the primewise bounds proves the lemma.
\end{proof}

\subsection{Prime assignments, factorials and scale constraints}

We have bounded a record by a product of token weights and reciprocal
prime values.  We now sum these products.  The exact subtree identity
will cancel their background exponential.  We must retain the factorials
for equal token types in order for this cancellation to remain exact.

For a fixed topology, record the tagged tokens in an ordered list and
record all list metadata: lengths, signs, attachments, suffix choices,
band choices, prime-slot sizes and equalities.  Apart from the token
shapes and prime values themselves, this costs
\begin{equation}\label{eq:trace-exceptional-entropy}
 \exp(O(N_{\mathrm{exc}}\log B)).
\end{equation}
Indeed there are at most $N_{\mathrm{exc}}$ objects and slots; each
positional index has at most a fixed power of $B$ choices, and their
equality pattern has at most
$N_{\mathrm{exc}}^{N_{\mathrm{exc}}}$ possibilities.  Multiple tagged
tokens may share a prime, and each resulting group has one prime
variable.  None of these variables is shared with an untagged token.

For the untagged tokens specify multiplicities $j_{a,Q}\ge0$ for their
types $(a,Q)$.  The actual prime values of one type form an unordered
set of $j_{a,Q}$ distinct primes.  Replace it temporarily by named
slots $1,\ldots,j_{a,Q}$ and divide the sum by
\begin{equation}\label{eq:trace-factorial}
 \prod_{a,Q}j_{a,Q}!.
\end{equation}
Every valid prime assignment has exactly this number of named
realizations.  Permuting values within one type preserves every tree
label, hence all return labels, the numerical line and the list.
In particular every bin condition is preserved.  After obtaining this
identity, we may enlarge the nonnegative sum by dropping distinctness
and other compatibility conditions.  Such an enlargement is a formal
sum of the established numerical majorants; it is not an application
of \Cref{lem:trace-integration} to new lines with coincident prime slots.

\begin{lemma}\label{lem:trace-pivots}
Fix the topology, token shapes, tagged/list groupings and untagged
multiplicities.  For every non-skipped good origin $u\notin\mathcal U$
choose deterministically one named core slot whose token has top $u$.
When summing prime assignments, the bin conditions on these edges give
an extra factor $O(B^{-1+\eta})$ for each chosen slot, while retaining
the unrestricted harmonic factor $v_a$ for every prime variable.
Together with the charge exponent and its correction in
\eqref{eq:trace-integration}, these savings are bounded by
\begin{equation}\label{eq:trace-pivot-saving}
 B^{-100r_*}\exp(O(\ell))
       B^{-(1-\eta)(|\mathcal G|-N_{\mathrm{exc}})}.
\end{equation}
\end{lemma}

\begin{proof}
Choose one eligible type at each origin using any fixed order, and
then choose its first named slot.  Distinct origins choose different
slots.  Because $s(u)=1$, a nontrivial subtree with top $u$ contains
the edge leaving $u$.

Fix all other prime values and order the selected slots ancestor-first
by their tops.  If the token of a selected slot occurs on the edge
leaving $u$, its top is $u$ or an ancestor of $u$.  Thus that edge's
product contains its selected prime and only earlier selected primes.
The bin requirement has the form
\[
 H<p_i D_i(p_1,\ldots,p_{i-1})\le\tau H,
\]
where $D_i$ is positive and all its other factors have been fixed.
The harmonic sum of a core prime in such an interval is
$O(B^{-1+\eta})$, uniformly in its position.  To see this, write the
interval as $(x,\tau x]$.  If it meets the core band, then
$x\ge \exp(B^{1-\eta})/\tau$, and the prime-counting upper bound gives
\[
 \sum_{\substack{x<p\le\tau x\\p\in\mathcal C}}\frac1p
 \le\frac{\pi(\tau x)}x
 \ll_\tau\frac1{\log x}\ll B^{-1+\eta}.
\]
Eliminate the selected variables in \emph{reverse} ancestor order.
Each bound is uniform in the remaining earlier values, so induction
gives one such factor per selected slot.  Since $v_C\ge1$ for large
$B$, each bound is also $O(B^{-1+\eta})v_C$.  All unselected variables
keep their unrestricted harmonic sums.  This choice of named pivots
requires no enumeration of possible pivots, and the bin condition
holds for every slot permutation counted in \eqref{eq:trace-factorial}.

The charge contribution is
\[
 -\kappa(v_C-T\sqrt{v_C})r_*
 -\bigl[\kappa(v_C-T\sqrt{v_C})
           -v_C\beta_C(e^\kappa-1)\bigr]|\mathcal U|.
\]
Both bracketed coefficients exceed $100\log B$ for large $B$:
$v_C=(\eta+o(1))\log B$, $\kappa\eta=400$, and
$\beta_C(e^\kappa-1)<1$.  Every non-skipped good edge therefore pays
either its pivot saving or the stronger factor $B^{-100}$ from its
origin in $\mathcal U$.  At most $N_{\mathrm{exc}}$ good edges were
skipped.  Absorbing the constants for at most $\ell$ pivots gives
\eqref{eq:trace-pivot-saving}.
\end{proof}

\begin{lemma}\label{lem:trace-summation}
For sufficiently large $B$, the total conditional main term supplied
by \Cref{prop:cylinder-reduction}, including its
$\exp(O(M_0\log B))$ prefactor, is at most
$B^{-(1+2c_*)\ell}$.
\end{lemma}

\begin{proof}
First apply \Cref{lem:trace-integration,lem:trace-pivots} to valid
assigned records.  The pivot bound is uniform in the detailed record
and may be factored out for each topology.

For the tagged/list groups, unrestricted prime sums cost at most
$(1+v_C+v_Z)^{O(N_{\mathrm{exc}})}$.  A tagged subtree shape can be
summed with its weight using \eqref{eq:subtree-identity}, at a cost
at most $\ell+1$; a ghost has at most $\ell$ possible edges and weight
$2$.  Ignoring compatibility between token shapes increases these
positive sums.  Thus all these shape costs are
$(O(\ell+1))^{O(N_{\mathrm{exc}})}$.  Together with
\eqref{eq:trace-exceptional-entropy}, the lost skipped-edge savings
and the cylinder prefactor, their logarithm is
\begin{equation}\label{eq:trace-small-costs}
 O(N_{\mathrm{exc}}\log B+M_0\log B)=o(\ell).
\end{equation}

For the untagged tokens, retain \eqref{eq:trace-factorial} and sum
over unrestricted multiplicities.  Their total is at most
\begin{align*}
 \prod_{a,Q}\sum_{j\ge0}
       \frac{(v_a\widetilde u_a(Q))^j}{j!}
 &=\exp\left\{\sum_a v_a\sum_Q\widetilde u_a(Q)\right\}.
\end{align*}
These are finite products of positive convergent series.  They cancel
the unmodified background terms
$-v_CS_C(\mathscr T)-v_ZS_Z(\mathscr T)$ in
\eqref{eq:trace-integration}.  By
\eqref{eq:modified-subtree-sum}, the remaining exponential is
\[
 \exp\left\{-\frac\theta2v_Z|\mathcal G|
              +\theta v_Z|E(\mathscr T)\setminus\mathcal G|\right\}
 \le B^{-(\eps/10)|\mathcal G|
                +\eps|E(\mathscr T)\setminus\mathcal G|}
\]
for large $B$, since
$(\theta/2)(\eps-\eta)>\eps/10$ and
$\theta(\eps-\eta)<\eps$.

It remains to check the numerical margin including the topology count.
Write $g=|\mathcal G|$, $b=|E(\mathscr T)\setminus\mathcal G|$ and
$r=r_*$.  Then $\ell=g+b+r$ and $b\le10(r+1)$.
Set $\delta=\eps/10-\eta=9\eta$.  Combining the preceding bounds,
\eqref{eq:trace-pivot-saving} and the count in
\Cref{lem:trace-topology}, the negative base-$B$ exponent is at least
\[
 (1+\delta)g+(99-o(1))r-\eps b-o(\ell).
\]
Here $4^\ell$ and the fixed constants per pivot contribute
$O(\ell/\log B)=o(\ell)$, and
$(\ell+1)^r=B^{(1+o(1))r}$.  Substitution of $g=\ell-b-r$ and
$b\le10(r+1)$ gives the lower bound
\[
 (1+\delta)\ell
 +(88-11\delta-10\eps-o(1))r
 -10(1+\delta+\eps)-o(\ell).
\]
The coefficient of $r$ is positive, and
$\delta-2c_*=7\eta>0$.  The slack absorbs the fixed last term, the
displayed $o(\ell)$ terms, and the sum over at most $\ell+1$ possible
values of $r$.  This proves the claimed bound.
\end{proof}

\begin{proof}[Completion of the proof of \Cref{thm:high-trace}]
The density assertion is \Cref{lem:bad-density}.  Apply
\Cref{prop:retained-lines} to split the trace sum into retained and
discarded configurations.  For the retained part, the triangle
inequality bounds the absolute full expectation by
$\E_{\mathcal F}R_{\boldsymbol d}
 |\E_{\mathcal S\setminus\mathcal F}
 K_{\boldsymbol d}\prod_j\1_Y(n+b_j)|$.
Now apply \Cref{prop:cylinder-reduction}.  Their discarded
contributions are $O(\exp(-B^{1+\eps/2}))$.  The retained conditional
main term is bounded by \Cref{lem:trace-summation}.  Since
$\ell=2\lfloor B\rfloor$,
\[
 B^{-(1+2c_*)\ell}+O(\exp(-B^{1+\eps/2}))
 \le B^{-(1+c_*)\ell}
\]
for sufficiently large $B$.  This proves the trace assertion and
completes the theorem.
\end{proof}
